% GENERATED FILE. Edit canonical lessons, then run npm run generate:books.
% canonical-source-hash: fnv1a64:74a54a8ec908dd46
% canonical-lessons: JA-C368-zentai, JA-C368-bubun, JA-C368-ryouhou, JA-C368-katahou, JA-C368-katayoru

\chapter{All or Part}
\label{ch:ja-all-or-part}
\hlchaptermodality{368}

\begin{chapteropening}
\textbf{By the end of this chapter:} \emph{I can talk about the whole of something or one part of it.}

Complete the last lesson of chapter 368: I can talk about the whole of something or one part of it.
\end{chapteropening}

\section[\texorpdfstring{\ja{ぜんたい} (zentai)}{ぜんたい (zentai)}]{\ja{ぜんたい} (zentai) — the whole}

\label{lesson:JA-C368-zentai}



\begin{quote}
Before the new one: say the Japanese for a degree, then the Japanese for a proportion.
\end{quote}

\subsection*{You'll want to know: \ja{ぜんたい}}
\textbf{\ja{ぜんたい}} — \emph{zentai} — \textquotedblleft{}the whole\textquotedblright{}.

\begin{quote}
\textbf{\ja{ぜんたいを} \ja{よんだら、いけんを} \ja{言}\ja{って} \ja{ください。}} — \emph{zentai o yondara, iken o itte kudasai} — \textquotedblleft{}When you have read the whole thing, please tell me what you think.\textquotedblright{}
\end{quote}

\begin{cousinweb}
\textbf{\ja{ぜんたい}} (\emph{zentai}) is Sino-Japanese (kango): two Chinese characters read with their on'yomi, readings borrowed from Chinese. They are \emph{zen}, \textquotedblleft{}all\textquotedblright{} (Middle Chinese \emph{dzjwen}) and \emph{tai}, \textquotedblleft{}body\textquotedblright{} (\emph{thejX}) (KANJIDIC2; Middle Chinese in Baxter's transcription).
\end{cousinweb}

\begin{grammarlens}[title={What you've built}]
One more word for saying what happens if.
\end{grammarlens}

\subsection*{Your turn}
\begin{itemize}
\raggedright
  \item \emph{Say it:} \emph{zentai}
  \item \emph{Say it:} \emph{zentai}, once more
  \item \emph{Say it:} \emph{zentai o yondara, iken o itte kudasai} — When you have read the whole thing, please tell me what you think.
  \item \emph{Recall:} say the Japanese for a degree, then the Japanese for a proportion, then say \emph{zentai} again
\end{itemize}

\begin{tcolorbox}[breakable,colback=teal!4,colframe=teal!35!black,title={Before you move on}]
What does \ja{ぜんたい} mean? (\textquotedblleft{}The whole\textquotedblright{}.) Say it once more.
\end{tcolorbox}
\section[\texorpdfstring{\ja{ぶぶん} (bubun)}{ぶぶん (bubun)}]{\ja{ぶぶん} (bubun) — a part; a section}

\label{lesson:JA-C368-bubun}



\begin{quote}
Before the new one: say the Japanese for a proportion, then the Japanese for the whole.
\end{quote}

\subsection*{You'll want to know: \ja{ぶぶん}}
\textbf{\ja{ぶぶん}} — \emph{bubun} — \textquotedblleft{}a part; a section\textquotedblright{}.

\begin{quote}
\textbf{\ja{この} \ja{ぶぶんに} \ja{しつもんが} \ja{あったら、きいて} \ja{ください。}} — \emph{kono bubun ni shitsumon ga attara, kiite kudasai} — \textquotedblleft{}If you have a question about this part, please ask.\textquotedblright{}
\end{quote}

\begin{cousinweb}
\textbf{\ja{ぶぶん}} (\emph{bubun}) is Sino-Japanese (kango): two Chinese characters read with their on'yomi, readings borrowed from Chinese. They are \emph{bu}, \textquotedblleft{}section\textquotedblright{} (Middle Chinese \emph{buX}) and \emph{bun}, \textquotedblleft{}part\textquotedblright{} (\emph{pjun}) (KANJIDIC2; Middle Chinese in Baxter's transcription).
\end{cousinweb}

\begin{grammarlens}[title={What you've built}]
One more word for saying what happens if.
\end{grammarlens}

\subsection*{Your turn}
\begin{itemize}
\raggedright
  \item \emph{Say it:} \emph{bubun}
  \item \emph{Say it:} \emph{bubun}, once more
  \item \emph{Say it:} \emph{kono bubun ni shitsumon ga attara, kiite kudasai} — If you have a question about this part, please ask.
  \item \emph{Recall:} say the Japanese for a proportion, then the Japanese for the whole, then say \emph{bubun} again
\end{itemize}

\begin{tcolorbox}[breakable,colback=teal!4,colframe=teal!35!black,title={Before you move on}]
What does \ja{ぶぶん} mean? (\textquotedblleft{}A part; a section\textquotedblright{}.) Say it once more.
\end{tcolorbox}
\section[\texorpdfstring{\ja{りょうほう} (ryōhō)}{りょうほう (ryōhō)}]{\ja{りょうほう} (ryōhō) — both}

\label{lesson:JA-C368-ryouhou}



\begin{quote}
Before the new one: say the Japanese for the whole, then the Japanese for a section.
\end{quote}

\subsection*{You'll want to know: \ja{りょうほう}}
\textbf{\ja{りょうほう}} — \emph{ryōhō} — \textquotedblleft{}both\textquotedblright{}.

\begin{quote}
\textbf{\ja{りょうほう} \ja{かったら、たかいでしょう。}} — \emph{ryōhō kattara, takai deshō} — \textquotedblleft{}If you buy both, it will probably be expensive.\textquotedblright{}
\end{quote}

\begin{cousinweb}
\textbf{\ja{りょうほう}} (\emph{ryōhō}) is Sino-Japanese (kango): two Chinese characters read with their on'yomi, readings borrowed from Chinese. They are \emph{ryō}, \textquotedblleft{}both\textquotedblright{} (Middle Chinese \emph{ljangX}) and \emph{hō}, \textquotedblleft{}direction\textquotedblright{} (\emph{pjang}) (KANJIDIC2; Middle Chinese in Baxter's transcription).
\end{cousinweb}

\begin{grammarlens}[title={What you've built}]
One more word for saying what happens if.
\end{grammarlens}

\subsection*{Your turn}
\begin{itemize}
\raggedright
  \item \emph{Say it:} \emph{ryōhō}
  \item \emph{Say it:} \emph{ryōhō}, once more
  \item \emph{Say it:} \emph{ryōhō kattara, takai deshō} — If you buy both, it will probably be expensive.
  \item \emph{Recall:} say the Japanese for the whole, then the Japanese for a section, then say \emph{ryōhō} again
\end{itemize}

\begin{tcolorbox}[breakable,colback=teal!4,colframe=teal!35!black,title={Before you move on}]
What does \ja{りょうほう} mean? (\textquotedblleft{}Both\textquotedblright{}.) Say it once more.
\end{tcolorbox}
\section[\texorpdfstring{\ja{かたほう} (katahō)}{かたほう (katahō)}]{\ja{かたほう} (katahō) — one of a pair; one side}

\label{lesson:JA-C368-katahou}



\begin{quote}
Before the new one: say the Japanese for a section, then the Japanese for both.
\end{quote}

\subsection*{You'll want to know: \ja{かたほう}}
\textbf{\ja{かたほう}} — \emph{katahō} — \textquotedblleft{}one of a pair; one side\textquotedblright{}.

\begin{quote}
\textbf{\ja{てぶくろを} \ja{かたほう} \ja{なくしたら、こまるでしょう。}} — \emph{tebukuro o katahō nakushitara, komaru deshō} — \textquotedblleft{}If you lose one glove, you will be stuck.\textquotedblright{}
\end{quote}

\begin{cousinweb}
\textbf{\ja{かたほう}} (\emph{katahō}) is built from the native \emph{kata}, \textquotedblleft{}one of a pair\textquotedblright{}, a kun'yomi, and the Sino-Japanese \emph{hō}, \textquotedblleft{}direction\textquotedblright{}, an on'yomi borrowed from Chinese (Middle Chinese \emph{pjang}) (KANJIDIC2; JMdict; Middle Chinese in Baxter's transcription).
\end{cousinweb}

\begin{grammarlens}[title={What you've built}]
One more word for saying what happens if.
\end{grammarlens}

\subsection*{Your turn}
\begin{itemize}
\raggedright
  \item \emph{Say it:} \emph{katahō}
  \item \emph{Say it:} \emph{katahō}, once more
  \item \emph{Say it:} \emph{tebukuro o katahō nakushitara, komaru deshō} — If you lose one glove, you will be stuck.
  \item \emph{Recall:} say the Japanese for a section, then the Japanese for both, then say \emph{katahō} again
\end{itemize}

\begin{tcolorbox}[breakable,colback=teal!4,colframe=teal!35!black,title={Before you move on}]
What does \ja{かたほう} mean? (\textquotedblleft{}One of a pair; one side\textquotedblright{}.) Say it once more.
\end{tcolorbox}
\section[\texorpdfstring{\ja{かたよる} (katayoru)}{かたよる (katayoru)}]{\ja{かたよる} (katayoru) — to be one-sided; to lean to one side}

\label{lesson:JA-C368-katayoru}



\begin{quote}
Before the new one: say the Japanese for both, then the Japanese for one of a pair.
\end{quote}

\subsection*{You'll want to know: \ja{かたよる}}
\textbf{\ja{かたよる}} — \emph{katayoru} — \textquotedblleft{}to be one-sided; to lean to one side\textquotedblright{}.

\begin{quote}
\textbf{\ja{たべものが} \ja{かたよったら、びょうきに} \ja{なるでしょう。}} — \emph{tabemono ga katayottara, byōki ni naru deshō} — \textquotedblleft{}If your diet is unbalanced, you will probably fall ill.\textquotedblright{}
\end{quote}

Polite \textbf{\ja{かたよります}}, \emph{katayorimasu}; a godan verb. Past (polite): \textbf{\ja{いけんが} \ja{かたよりました。}}, \emph{iken ga katayorimashita}.

\begin{cousinweb}
\textbf{\ja{かたよる}} (\emph{katayoru}) is a native Japanese word (wago), inherited rather than borrowed from Chinese, and its older origin is not settled. Its kanji was attached to it later as a kun'yomi, a native reading given to a Chinese character (JMdict, KANJIDIC2).
\end{cousinweb}

\begin{grammarlens}[title={What you've built}]
One more word for saying what happens if.
\end{grammarlens}

\subsection*{Your turn}
\begin{itemize}
\raggedright
  \item \emph{Say it:} \emph{katayoru}
  \item \emph{Say it:} \emph{katayoru}, once more
  \item \emph{Say it:} \emph{tabemono ga katayottara, byōki ni naru deshō} — If your diet is unbalanced, you will probably fall ill.
  \item \emph{Recall:} say the Japanese for both, then the Japanese for one of a pair, then say \emph{katayoru} again
\end{itemize}

\begin{tcolorbox}[breakable,colback=teal!4,colframe=teal!35!black,title={Before you move on}]
What does \ja{かたよる} mean? (\textquotedblleft{}To be one-sided; to lean to one side\textquotedblright{}.) Say it once more.
\end{tcolorbox}
